\begin{table}[!ht]
\centering
\caption{Balanceamento entre trabalhadores brancos e negros \emph{antes} de qualquer
controle --- a tabela que mostra de que é feita a comparação (Angrist \& Pischke, cap.~3).
Diferença padronizada: $d = (\bar x_b - \bar x_n)/\sqrt{(s_b^2+s_n^2)/2}$; valores acima de
$0{,}10$ em módulo indicam desequilíbrio relevante. População completa da PEA com renda
positiva ($N = 7.694.095$; 3.259.472 brancos e 4.434.623 negros).}
\label{tab:balanceamento}
\small
\begin{tabular}{lrrrr}
\toprule
Variável & Brancos & Negros & Diferença & $d$ \\
\midrule
log do rendimento & 7,577 & 7,151 & 0,425 & 0,47 \\
Idade (centrada) & −1,638 & −3,110 & 1,472 & 0,11 \\
Sexo feminino & 0,437 & 0,399 & 0,038 & 0,08 \\
Fundamental completo & 0,003 & 0,009 & −0,006 & −0,07 \\
Médio completo & 0,046 & 0,041 & 0,006 & 0,03 \\
Superior completo & 0,198 & 0,278 & −0,080 & −0,19 \\
Pós-graduação & 0,006 & 0,009 & −0,004 & −0,04 \\
Escolaridade não registrada & 0,737 & 0,653 & 0,084 & 0,18 \\
log(horas trabalhadas) & 3,644 & 3,598 & 0,045 & 0,11 \\
Reside em área urbana & 0,813 & 0,773 & 0,039 & 0,10 \\
Emprego formal & 0,493 & 0,439 & 0,054 & 0,11 \\
Conta própria & 0,277 & 0,282 & −0,006 & −0,01 \\
Trabalho doméstico & 0,049 & 0,079 & −0,030 & −0,12 \\
CBO: dirigente & 0,055 & 0,023 & 0,032 & 0,16 \\
CBO: profissional & 0,149 & 0,079 & 0,071 & 0,22 \\
CBO: técnico & 0,090 & 0,066 & 0,023 & 0,09 \\
\% negro na UPA ($z$) & −0,701 & 0,331 & −1,032 & −1,16 \\
Desemprego na UPA ($z$) & −0,322 & 0,099 & −0,422 & −0,46 \\
Educ. média na UPA ($z$) & −0,134 & 0,162 & −0,297 & −0,30 \\
\bottomrule
\end{tabular}
\normalsize
\par\smallskip
\footnotesize\emph{Suporte comum:} 97,3\% das UPAs têm
trabalhadores dos dois grupos, e 100,0\% das células
UF~$\times$~escolaridade também (100,0\% das observações
estão em células com os dois grupos) --- a comparação não depende de extrapolar para
regiões da amostra onde só existe um dos grupos.
\end{table}
